\section{季度时间线：为什么这一季不是普通更新}

本期季报把 2026 年第一季度描述为一个“推背感”极强的阶段：模型从聊天进入 agentic 任务，Coding 让研发周期被压缩，前沿实验室的产品和 feature 发布频率显著上升。这里的关键不是某个单点模型，而是能力、工具、组织和工作流同时发生变化。

\subsection{从 2023 到 2026 的季报主线}

\noindent\begin{longtable}{p{0.20\textwidth}p{0.36\textwidth}p{0.35\textwidth}}
\toprule
阶段 & 主线 & 本期如何继承 \\
\midrule
2023 & ChatGPT 释放聊天范式，开源开始追赶 & 第一幕是“模型会说话、会回答”。 \\
2024 & AGI 基建、推理、self-play RL、o1 叙事出现 & 模型开始从回答转向推理和行动。 \\
2025 & 分化收敛、全家桶、垂直整合、AI War & 模型公司开始围绕产品和生态整合。 \\
2026 Q1 & Coding/Agent 成为第二幕主线 & 模型进入数字工作流，开始改造研发组织。 \\
\bottomrule
\end{longtable}

\begin{importantbox}{季度报告的价值}
单条新闻容易夸大短期噪音；季度视角能观察哪些判断持续成立。本期最重要的延续判断是：Coding 去年还是预言，今年已经变成海啸。
\end{importantbox}

\subsection{本章小结}

EP136 的意义在于把过去几季的判断串起来：模型从聊天、推理、产品生态，走到 Coding/Agent 这个能直接改变工作流的阶段。

